\section{Introduction}

\IEEEPARstart{T}{he} rapid advancement of connected vehicle technologies has transformed modern transportation infrastructure, with Vehicular Ad-hoc Networks (VANETs) serving as the backbone for next-generation intelligent transportation systems. By facilitating seamless communication between vehicles and roadside infrastructure, VANETs enable critical applications ranging from autonomous driving coordination to real-time hazard detection and cooperative traffic optimization \cite{engoulou2014vanet}. The computational demands of these applications continue to escalate, requiring sophisticated processing capabilities that exceed the limited onboard resources available in individual vehicles. Consequently, task offloading to edge computing infrastructure has become essential for meeting the stringent performance requirements of modern vehicular services \cite{buttazzo1997hard}. 

To address this computational resource limitation, fog computing has emerged as a promising paradigm that deploys processing capabilities at Roadside Units (RSUs) and edge servers positioned throughout the network \cite{mouradian2017comprehensive}. By bringing computational resources closer to vehicles, fog computing reduces communication latency compared to traditional cloud-based solutions while providing sufficient processing power for complex vehicular applications. However, a critical observation reveals that many computational tasks in VANETs exhibit significant redundancy. Multiple vehicles frequently request identical computations—such as optimal route calculations for common destinations, traffic density estimations for shared road segments, or environmental data processing from overlapping sensor coverage areas. Repeatedly executing these duplicate tasks consumes substantial computational resources and energy at fog nodes \cite{wu2021design}.

Task result caching offers an elegant solution by storing outputs of previously executed tasks and reusing them for subsequent identical requests. By preventing repetitive task execution, caching conserves both energy and computational resources at fog nodes, freeing these resources for executing other critical tasks and improving overall system throughput. When a vehicle submits a task request, the system first queries the cache; upon finding a cache hit, the stored result is immediately returned without re-execution, dramatically reducing response time, energy consumption, and computational overhead. However, implementing effective caching in VANETs introduces unique challenges. The constant mobility of vehicles creates highly dynamic access patterns where task popularity fluctuates based on geographic location, time of day, and traffic conditions \cite{domingos2016data}. Limited storage capacity at fog nodes demands intelligent cache management strategies, while the heterogeneous nature of vehicular tasks—characterized by diverse computational requirements, varying urgency levels, and different data volumes—necessitates sophisticated decision-making beyond simple frequency-based heuristics.

Traditional caching policies such as Least Recently Used (LRU), Least Frequently Used (LFU), and First-In-First-Out (FIFO) rely on static rules that cannot adapt to the unique dynamics of vehicular networks. Machine learning techniques present a promising alternative by enabling data-driven, adaptive caching strategies. Federated Learning (FL) enables distributed model training across multiple edge nodes while preserving data privacy, analyzing task characteristics to predict future popularity and proactively populate caches \cite{yu2015cooperative}. Complementing this, Reinforcement Learning (RL) offers a framework for optimizing cache eviction decisions through continuous interaction with the system environment, dynamically adjusting to real-time conditions and balancing multiple competing goals such as maximizing cache hit rates and minimizing task completion times. However, existing research predominantly explores FL and RL independently for VANET caching, overlooking the potential synergies from combining these complementary approaches. Furthermore, most prior work concentrates on content caching for data dissemination rather than task result caching for computational offloading, which possesses unique attributes including execution deadlines, resource consumption profiles, and priority classifications. The integration of Software-Defined Networking (SDN) with VANETs offers centralized visibility and programmable control, creating opportunities for coordinated caching decisions across distributed fog infrastructure \cite{misra2019soft}, yet leveraging SDN for intelligent, learning-based cache management in vehicular fog computing remains largely unexplored.

Addressing these gaps, we introduce \textbf{FL-RL-TRC}, a novel dual-phase intelligent caching framework for task result management in Software-Defined VANETs that synergistically combines Federated Learning for anticipatory cache population with Deep Reinforcement Learning for adaptive eviction control. Phase 1 employs an FL advisor that examines task patterns using a comprehensive 12-dimensional feature representation encompassing request frequencies, task priority levels, data payload sizes, locality indicators, vehicle velocity measurements, RSU computational loads, exponential popularity decay factors, and one-hot encoded task category identifiers, enabling accurate popularity forecasting and proactive cache pre-loading. Phase 2 deploys a Deep Q-Network (DQN) agent that continuously monitors system conditions through a 9-dimensional state vector and selects from five eviction strategies: no-operation, LRU eviction, random eviction, LFU eviction, or TTL shrinking. This dual-phase design leverages FL's strength in pattern recognition with RL's capability for dynamic optimization.

The key contributions of FL-RL-TRC are as follows:

\begin{enumerate}
    \item \textbf{Novel Two-Phase Architecture}: We propose the first hybrid caching framework that combines FL-based proactive pre-placement with DRL-based reactive eviction for task result caching in SDN-based VANETs. This separation of concerns allows each learning component to focus on its strengths—FL for long-term pattern recognition and RL for short-term dynamic optimization—while maintaining overall system coherence.
    
    \item \textbf{Rich Feature Engineering}: We develop comprehensive feature vectors derived from state-of-the-art VANET caching literature, incorporating mobility characteristics (vehicle speed, direction), network load (RSU utilization), temporal dynamics (exponential popularity decay with tunable parameter $\lambda$), and task-specific attributes (type, priority, size, locality). This multi-faceted representation enables more accurate popularity prediction and eviction decisions compared to simple frequency-based approaches.
    
    \item \textbf{Systematic Parameter Optimization}: We conduct extensive experiments to optimize the popularity decay parameter $\lambda$ across values \{0.01, 0.02, 0.05, 0.1, 0.2\}, identifying $\lambda = 0.2$ as optimal for VANET environments. This finding reveals that faster popularity decay (content becomes less relevant after approximately 5 seconds) better matches the highly dynamic nature of vehicular networks, providing valuable insights for future caching system design.
    
    \item \textbf{Superior Performance}: FL-RL-TRC achieves 97.0\% cache hit rate on synthetic VANET traces with 50 vehicles, 5 RSUs, and 200 timesteps, significantly outperforming static policies (LFU/LRU/Random: 88.1\%) by 8.9 percentage points, while maintaining 81.6\% task success rate and 7.31 ms average completion time.
    
    \item \textbf{Modular Design}: Our architecture enables independent optimization of FL and RL components. The FL advisor can be trained offline and deployed via Federated Averaging (FedAvg), while the RL agent continuously learns online through experience replay and $\epsilon$-greedy exploration.
\end{enumerate}

We validate FL-RL-TRC through extensive experimentation on synthetically generated VANET traces with varying task characteristics including four task categories, multiple priority tiers, heterogeneous data sizes, and hard deadline constraints. The evaluation includes systematic parameter sweeps, ablation studies, and comprehensive comparisons against baseline policies. Results demonstrate that FL-RL-TRC substantially exceeds all static policies, with the RL agent exhibiting clear learning progression from 87.8\% hit rate at 50 timesteps to 97.0\% at 200 timesteps. Our parameter optimization identifies $\lambda = 0.2$ as optimal, indicating that task relevance diminishes rapidly (approximately 5-second half-life) in highly mobile vehicular environments.

The remainder of this paper proceeds as follows: Section II surveys related literature on vehicular caching, federated learning, and reinforcement learning. Section III formalizes the system model including SDN-VANET architecture, task model, and cache model. Section IV details FL-RL-TRC design including the FL advisor's 12-dimensional feature extraction and the DRL agent's 9-dimensional state representation. Section V describes implementation specifics including trace generation, popularity decay modeling, and training procedures. Section VI presents experimental evaluation with parameter optimization, comparative analysis, and ablation studies. Section VII concludes with key findings and future research directions.
